% ============================================================
% citas.tex — DOCUMENTO MÍNIMO DE LA PASADA DE tex4ht (SC-57)
% LO ESCRIBE preparar_citas_libro.sh: NO SE EDITA A MANO
% ============================================================
\documentclass{book}
% ============================================================
% citas-proyecto.tex — INSUMO DE LA PASADA DE CITAS (SC-57)
% ------------------------------------------------------------
% ARCHIVO DERIVADO — LO ESCRIBE gbpublisher ANTES DE HTML Y EPUB.
% SON LAS MISMAS DEFINICIONES DE latex/preambulo-proyecto.tex, DE
% LA MISMA FUNCIÓN. NO EDITAR A MANO: LOS CAMBIOS SE HACEN EN LA BASE.
% ============================================================

% --- IDIOMA ---
\def\gbidiomas{main=spanish}
\def\gbidiomaprincipal{spanish}
\newif\ifgbcjk   \gbcjkfalse

% --- BIBLIOGRAFÍA ---
\def\estandar{verona-modern}
\newif\ifgbbibporcapitulo   \gbbibporcapitulotrue
\def\gbarchivobib{./referencias/ref-l-DUEK.bib}


\usepackage{fontspec}
% EL MISMO IDIOMA Y LAS MISMAS COMILLAS QUE EL PDF
% (preambulo-contrato.tex, SECCIÓN IDIOMA)
\usepackage[\gbidiomas]{babel}
\usepackage[autostyle=true]{csquotes}
\input{cita-\estandar-config.tex}
\input{cita-\estandar-html.tex}
% EL .bib SE COPIA A LA CARPETA: \gbarchivobib ES RELATIVO A LA
% RAÍZ DEL PROYECTO, Y ESTE DOCUMENTO SE COMPILA EN OTRO LUGAR
\addbibresource{ref-l-DUEK.bib}
\usepackage{hyperref}
% EN UNA COMPILACIÓN COMÚN SOLO ENVUELVE; cita-tex4ht.cfg LA
% REDEFINE COMO <span id="gbcita-...">
\newcommand*{\gbcita}[2]{#2}
% LA MARCA DE CADA BIBLIOGRAFÍA: EN UNA COMPILACIÓN COMÚN NO
% ESCRIBE NADA; cita-tex4ht.cfg LA REDEFINE COMO <span id="gbbib-...">
\newcommand*{\gbbibliografia}[1]{}
\begin{document}

\begin{refsection}
\chapter*{fm-10}
\gbcita{fm-10-1}{\protect\parencite*{4181-DALLE2022}}
\gbcita{fm-10-2}{\protect\parencite[p. 15]{4182-HEREDIA2022}}
\gbcita{fm-10-3}{\protect\parencite[p. 172]{4178-ALVAREZ2022}}
\gbcita{fm-10-4}{\protect\parencite*{4183-MARCUS2023}}
\gbcita{fm-10-5}{\protect\parencite*{4177-BALSA2023}}
\gbcita{fm-10-6}{\protect\parencite*{4188-SEGURA2024}}
\gbcita{fm-10-7}{\protect\parencite[p. 20]{4166-HELLER1977}}
\gbcita{fm-10-8}{\protect\parencite*{4184-GUBERN1987}}
\gbcita{fm-10-9}{\protect\parencite{4179-CASTELLS1998}}
\gbcita{fm-10-10}{\protect\parencite*{4187-MITCHELSTEIN2022}}
\gbcita{fm-10-11}{\protect\parencite*[p. 22]{4175-ARFUCH2002}}
\gbcita{fm-10-12}{\protect\parencite*[p. 1]{4176-AVERETT2021}}
\gbcita{fm-10-13}{\protect\parencite[p. 25]{4186-MECCIA2022}}
\gbcita{fm-10-14}{\protect\parencite{4167-RAPLEY2014}}
\gbcita{fm-10-15}{\protect\parencite{4180-CHARMAZ2006}}
\gbcita{fm-10-16}{\protect\parencite{4302-SPOLITA2022}}
\gbcita{fm-10-17}{\protect\parencite{4185-CLEMENCEAU2016}}
\gbbibliografia{fm-10}
\gbImprimirBibliografia
\end{refsection}


\begin{refsection}
\chapter*{a-01}
\gbcita{a-01-1}{\protect\parencite{4196-CHAVES2009}}
\gbcita{a-01-2}{\protect\parencite*{4226-ANDURRESTI2008}}
\gbcita{a-01-3}{\protect\parencite{4196-CHAVES2009}}
\gbcita{a-01-4}{\protect\parencite*{4227-REGUILLO2000}}
\gbcita{a-01-5}{\protect\parencite*{4195-MARGULIS2001}}
\gbcita{a-01-6}{\protect\parencite{4266-DESINGLY2000}}
\gbcita{a-01-7}{\protect\parencite*{4265-UNICEF2020}}
\gbcita{a-01-8}{\protect\parencite{4234-SAP2020}}
\gbcita{a-01-9}{\protect\parencite*{4198-FAILLA2021}}
\gbcita{a-01-10}{\protect\parencite*{4191-CARDINAUX2020}}
\gbcita{a-01-11}{\protect\parencite{4197-DUSSEL2020}}
\gbcita{a-01-12}{\protect\parencite{4233-MOGUILLANSKY2022}}
\gbcita{a-01-13}{\protect\parencite*{4197-DUSSEL2020}}
\gbcita{a-01-14}{\protect\parencite[p. 5]{4197-DUSSEL2020}}
\gbcita{a-01-15}{\protect\parencite*{4201-GIDDENS1997}}
\gbcita{a-01-16}{\protect\parencite{4230-SIMONI2022}}
\gbcita{a-01-17}{\protect\parencite{4230-SIMONI2022}}
\gbcita{a-01-18}{\protect\parencite{4200-HAIDT2024}}
\gbcita{a-01-19}{\protect\parencite{4232-TUNON2022}}
\gbcita{a-01-20}{\protect\parencite*{4274-BARCALA2022}}
\gbcita{a-01-21}{\protect\parencite{4233-MOGUILLANSKY2022}}
\gbcita{a-01-22}{\protect\parencite*{4231-TUNON2021}}
\gbcita{a-01-23}{\protect\parencite{4199-GARDINER1997}}
\gbcita{a-01-24}{\protect\parencite{4228-HOCHSCHILD2008,4229-PAUTASSI2013}}
\gbcita{a-01-26}{\protect\parencite*[pp. 6-7]{4190-CAGGIANO2021}}
\gbcita{a-01-27}{\protect\parencite{4233-MOGUILLANSKY2022}}
\gbcita{a-01-28}{\protect\parencite{4189-ACEVEDO2021,4192-ANDERETESCHWAL2020,4182-HEREDIA2022}}
\gbbibliografia{a-01}
\gbImprimirBibliografia
\end{refsection}


\begin{refsection}
\chapter*{a-02}
\footnote{\gbcita{a-02-1}{\protect\parencite[p. 47]{4239-MARTIN2023}}}
\footnote{\gbcita{a-02-2}{\protect\parencites[p. 17]{4243-HODGES2020}[en][p. 199]{4250-MACCHIAROLA2021}}}
\gbcita{a-02-4}{\protect\parencite{4247-CHAVES2005}}
\gbcita{a-02-5}{\protect\parencites{4195-MARGULIS2001}[entre otros]{4227-REGUILLO2000}}
\gbcita{a-02-7}{\protect\parencite[p. 37]{4248-CHAVES2010}}
\gbcita{a-02-8}{\protect\parencite[p. 10]{4238-NAVARRO2022}}
\gbcita{a-02-9}{\protect\parencite{4238-NAVARRO2022,4251-KAPLAN2021}}
\gbcita{a-02-11}{\protect\parencite{4251-KAPLAN2021}}
\gbcita{a-02-12}{\protect\parencite{4303-DINAPOLI2022}}
\gbcita{a-02-13}{\protect\parencite{4241-MARTUCCELLI2021}}
\gbcita{a-02-14}{\protect\parencite{4249-DALLE2022}}
\gbcita{a-02-15}{\protect\parencite{4246-CHAVEZMOLINA2022}}
\gbcita{a-02-16}{\protect\parencite{4185-CLEMENCEAU2016}}
\gbcita{a-02-17}{\protect\parencite{4235-PEREZMARTIRENA2022}}
\gbcita{a-02-18}{\protect\parencite{4242-MARTUCCELLI2012}}
\gbcita{a-02-19}{\protect\parencite{4235-PEREZMARTIRENA2022}}
\gbcita{a-02-20}{\protect\parencite{4185-CLEMENCEAU2016}}
\gbcita{a-02-21}{\protect\parencite{4194-BRACCHI2016}}
\gbcita{a-02-22}{\protect\parencite{4236-SIMONETTI2024}}
\gbcita{a-02-23}{\protect\parencite{4185-CLEMENCEAU2016}}
\gbcita{a-02-24}{\protect\parencite{4245-LEMUS2024}}
\gbcita{a-02-25}{\protect\parencites{4237-SANTOSSHARPE2022}{4240-MORAN2021}{4244-KRICHESKY2021}[entre otros]{4245-LEMUS2024}}
\gbcita{a-02-29}{\protect\parencite[p. 58]{4240-MORAN2021}}
\gbcita{a-02-30}{\protect\parencite{4244-KRICHESKY2021}}
\gbcita{a-02-31}{\protect\parencite{4237-SANTOSSHARPE2022}}
\gbcita{a-02-32}{\protect\parencite{4238-NAVARRO2022}}
\gbbibliografia{a-02}
\gbImprimirBibliografia
\end{refsection}


\begin{refsection}
\chapter*{a-03}
\gbcita{a-03-1}{\protect\parencite[p. 91]{4246-CHAVEZMOLINA2022}}
\gbcita{a-03-2}{\protect\parencite{4255-SAUTU2022}}
\gbcita{a-03-3}{\protect\parencite{4254-ORTALE2022}}
\gbcita{a-03-4}{\protect\parencite[pp. 45-46]{4252-BENZA2022}}
\gbcita{a-03-5}{\protect\parencite[p. 17]{4253-BENZA2021}}
\gbcita{a-03-6}{\protect\parencite[pp. 14-15]{4252-BENZA2022}}
\gbcita{a-03-7}{\protect\parencite[p. 18]{4252-BENZA2022}}
\gbcita{a-03-8}{\protect\parencite[p. 25]{4253-BENZA2021}}
\gbcita{a-03-9}{\protect\parencite[p. 7]{4252-BENZA2022}}
\gbcita{a-03-10}{\protect\parencite{4182-HEREDIA2022}}
\gbbibliografia{a-03}
\gbImprimirBibliografia
\end{refsection}


\begin{refsection}
\chapter*{a-04}
\gbcita{a-04-1}{\protect\parencite*{4222-SIBILIA2008}}
\gbcita{a-04-2}{\protect\parencite*{4224-CANEVARO2020}}
\gbcita{a-04-3}{\protect\parencite{4262-DUEK2020}}
\gbcita{a-04-4}{\protect\parencite*{4261-ALMADA2023}}
\gbcita{a-04-5}{\protect\parencite*{4182-HEREDIA2022}}
\gbcita{a-04-6}{\protect\parencite*{4184-GUBERN1987}}
\gbcita{a-04-7}{\protect\parencite*{4256-SEGURA2022}}
\gbcita{a-04-8}{\protect\parencite*{4225-HARVEY1977}}
\gbcita{a-04-9}{\protect\parencite*{4223-ARZA2020}}
\gbcita{a-04-10}{\protect\parencite*{4221-WAINERMAN2003}}
\gbcita{a-04-11}{\protect\parencite{4263-BATTHYANY2020}}
\gbcita{a-04-12}{\protect\parencite{4257-MENESESMEDINA2023,4258-PEREYRA2023}}
\gbcita{a-04-14}{\protect\parencite{4210-MOGUILLANSKY2024}}
\gbcita{a-04-15}{\protect\parencite*{4259-SALVIA2021}}
\gbcita{a-04-16}{\protect\parencite*{4260-JACOVKIS2021}}
\gbcita{a-04-17}{\protect\parencite*{4275-CIPPEC2021}}
\gbbibliografia{a-04}
\gbImprimirBibliografia
\end{refsection}


\begin{refsection}
\chapter*{a-05}
\gbcita{a-05-1}{\protect\parencite{4219-DALLE2021}}
\gbcita{a-05-2}{\protect\parencite{4182-HEREDIA2022}}
\gbcita{a-05-3}{\protect\parencite{4218-CASTEL1995}}
\gbcita{a-05-4}{\protect\parencite{4220-SALVIA2018}}
\gbcita{a-05-5}{\protect\parencite*{4217-BECK1998}}
\gbcita{a-05-6}{\protect\parencite[p. 13]{4218-CASTEL1995}}
\gbbibliografia{a-05}
\gbImprimirBibliografia
\end{refsection}


\begin{refsection}
\chapter*{a-06}
\gbcita{a-06-1}{\protect\parencite{4205-HEID2021}}
\gbcita{a-06-2}{\protect\parencite{4211-CIRINO2022}}
\gbcita{a-06-3}{\protect\parencite{4215-MINISTERIO2021}}
\gbcita{a-06-4}{\protect\parencite{4216-CEPAL2020}}
\gbcita{a-06-5}{\protect\parencite{4205-HEID2021,4211-CIRINO2022}}
\gbcita{a-06-7}{\protect\parencite{4168-CUNZOLO2021}}
\gbcita{a-06-8}{\protect\parencite{4206-VENTURIELLO2023}}
\gbcita{a-06-9}{\protect\parencite{4207-AGUIRRECUNS2018}}
\gbcita{a-06-10}{\protect\parencite{4169-CERRI2016}}
\gbcita{a-06-11}{\protect\parencite{4169-CERRI2016}}
\gbcita{a-06-12}{\protect\parencite{4203-SMITH2007}}
\gbcita{a-06-13}{\protect\parencite[p. 146]{4207-AGUIRRECUNS2018}}
\gbcita{a-06-14}{\protect\parencite{4174-BATTHYANY2021}}
\gbcita{a-06-15}{\protect\parencite{4204-VEGASOLIS2017}}
\gbcita{a-06-16}{\protect\parencite{4174-BATTHYANY2021}}
\gbcita{a-06-17}{\protect\parencite{4212-CEMINARI2018}}
\gbcita{a-06-18}{\protect\parencite[p. 28]{4213-RICO2016}}
\gbcita{a-06-19}{\protect\parencite{4202-BONIOLO2024,4204-VEGASOLIS2017,4174-BATTHYANY2021}}
\gbcita{a-06-22}{\protect\parencite[p. 27]{4173-LOPEZ2009}}
\gbcita{a-06-23}{\protect\parencite{4202-BONIOLO2024,4210-MOGUILLANSKY2024}}
\gbcita{a-06-25}{\protect\parencite{4174-BATTHYANY2021}}
\gbcita{a-06-26}{\protect\parencite*{4202-BONIOLO2024}}
\gbcita{a-06-27}{\protect\parencite{4206-VENTURIELLO2023}}
\footnote{\gbcita{a-06-28}{\protect\parencite{4170-ALGRANTI2010}}}
\footnote{\gbcita{a-06-29}{\protect\parencite{4171-BORRELLI2022}}}
\gbcita{a-06-30}{\protect\parencite{4172-IACUB2017}}
\gbcita{a-06-31}{\protect\parencite{4208-ARBER2003}}
\gbcita{a-06-32}{\protect\parencite{4214-GONZALVEZTORRALBO2021}}
\gbcita{a-06-33}{\protect\parencite{4209-BOYD2007}}
\gbcita{a-06-34}{\protect\parencite{4202-BONIOLO2024,4210-MOGUILLANSKY2024}}
\gbbibliografia{a-06}
\gbImprimirBibliografia
\end{refsection}

\chapter*{bm-10}

\chapter*{bm-11}

% BIBLIOGRAFÍA CONSOLIDADA: UNA SOLA, AL FINAL, COMO EN EL PDF.
% POR CAPÍTULO LA CIERRA CADA PIEZA EN SU refsection (GV-102)
\ifgbbibporcapitulo\else\gbbibliografia{libro}\gbImprimirBibliografia\fi
\end{document}
